\documentclass[11pt,a4paper]{article}
\usepackage[T1]{fontenc}
\usepackage{lmodern}
\usepackage{amsmath,amssymb,amsthm,mathtools}
\usepackage[margin=27mm,headheight=14pt]{geometry}
\usepackage{microtype}
\usepackage{longtable,booktabs,array}
\usepackage{enumitem}
\usepackage{xurl}
\usepackage[hidelinks,unicode,hypertexnames=false]{hyperref}
\usepackage{fancyhdr}
\pagestyle{fancy}
\fancyhf{}
\fancyhead[L]{\small The Calkin Algebra as a Test for Nuclearity}
\fancyhead[R]{\small\thepage}
\renewcommand{\headrulewidth}{0.3pt}
\setlength{\parindent}{1.2em}
\setlength{\parskip}{0pt}
\setlength{\emergencystretch}{2em}
\allowdisplaybreaks[2]
\hypersetup{pdftitle={The Calkin Algebra as a Test for Nuclearity},pdfauthor={Hikaru Awazu}}
\title{The Calkin Algebra as a Test for Nuclearity}
\author{Hikaru Awazu}
\date{}
\begin{document}
\maketitle
\begin{abstract}
Let \(Q=\mathcal B(\ell^2)/\mathcal K(\ell^2)\) be the Calkin algebra on a fixed separable infinite-dimensional Hilbert space.
This manuscript proposes a proof that
a complex \(C^*\)-algebra \(B\) is nuclear precisely when the canonical map
\(B\otimes_{\max}Q\longrightarrow B\otimes_{\min}Q\) is injective.
No unitality, separability, \(\sigma\)-unitality or initial exactness of \(B\) is assumed.
The argument first extracts a finite unitary family and one normal state in a literal corner of \(B^{**}\).
A cyclic-lamp construction then gives a single completely positive moment map defined on the whole Calkin algebra.
A Hilbert--Schmidt frame estimate supplies its normal source extension.
Finally finite normal tests select coefficients in \(B\) itself, and exactness supplies one ideal cut preserving a fixed state.
The proposed positive kernel element has state value greater than \(27/256\) and norm at most \(1/8\).
\end{abstract}

\begin{quote}\small
\textbf{Provenance and verification status.}
This manuscript was generated by GPT-6 and is presented by Hikaru Awazu as an unreviewed proof candidate.
The Lean formalization of the main theorem and its dependency closure has been rebuilt from source.
The correspondence between the formalization and this written proof, and independent human review, are not yet complete.
The established inputs are listed separately in Section 1.
All mathematical dependencies used in the argument are either stated standard inputs or supplied in the manuscript.
\end{quote}

\tableofcontents
\setcounter{section}{-1}
\section{Statement, conventions and notation}
Fix \(H_*=\ell^2(\mathbb N)\) and \(Q_*=\mathcal B(H_*)/\mathcal K(H_*)\).
For any complex \(C^*\)-algebra \(B\), let
\[
\chi_B:B\otimes_{\max}Q_*\longrightarrow B\otimes_{\min}Q_*
\]
be the canonical surjective \(*\)-homomorphism extending the identity on \(B\odot Q_*\).
The symbol \(\odot\) denotes the algebraic tensor product.

\paragraph{Main theorem (proposed).} For every such \(B\),
\begin{equation}\tag{T}
\chi_B\text{ is injective if and only if }B\text{ is nuclear}.
\end{equation}
This is the fixed-Calkin tensor test considered here. Injectivity means equality of the two completed \(C^*\)-norms, since an injective \(*\)-homomorphism is isometric. It does not require equality of norms in every individual representation. The implication from nuclearity follows from the definition of nuclearity. The case \(B=0\) is immediate. We prove the other implication for arbitrary nonzero \(B\), using Section~1 to obtain exactness.

All Hilbert-space inner products are linear in the first variable.
The conjugate \(C^*\)-algebra \(\overline N\) has conjugated scalar multiplication; it is not being identified silently with the opposite algebra. In a normal standard form of \(N\) on \(H_N\), \(L_N\) is the faithful left action, \(J_N\) the conjugation, and
\[
r_N(x)=J_N L_N(x^*)J_N,\qquad r_N(x)r_N(y)=r_N(yx).
\]
For a normal state \(\phi\), \(\xi_\phi\) is its unit vector in the natural cone.
A UCP map is unital and completely positive; normality of a positive map between von Neumann algebras means preservation of bounded increasing positive nets, equivalently ultraweak continuity. A trace called n.s.f.\ is normal, semifinite and faithful. \(N_*\) denotes the normal predual, whereas \(N^*\) denotes the full Banach dual. Corners have their relative units. Indexed families retain repetitions. All strong/ultraweak limits may be nets unless a sequence is specified. No target Hilbert space or predual is assumed separable.

\paragraph{Notation.}
\begin{description}
\item[\(B,V=B^{**},\kappa\)] Original \(C^*\)-algebra, canonical bidual, inclusion.
\item[\(H_*,Q_*,\chi_B\)] Fixed separable Hilbert space, Calkin algebra, canonical map.
\item[\(M=p_M V p_M\)] \(\sigma\)-finite noninjective type part in Section~2.
\item[\(C_0,C_2,j_0,j_1\)] Modular core, double product, canonical inclusions.
\item[\(E_0,E_1,\Xi\)] Two bounded returns and normal onto isomorphism \(C_2\to M\).
\item[\(F_0,F,i,E,\tau\)] Finite algebras, embedding, return, finite trace state.
\item[\(e,N=eVe\)] Final literal target corner and its relative unit \(e\).
\item[\(\chi,\psi,\phi\)] Faithful modular state, invariant witness, normal state.
\item[\(v_j,q,r\)] Target unitaries, paired min norm, cone correlation.
\item[\(G=\mathbb F(s,t),\Gamma\)] Free group and cyclic-lamp semidirect-product group.
\item[\(\Omega,R\)] Uniform lamp spectrum and \(L^\infty(\Omega)\).
\item[\(C_f,C,D\)] \(C^*(\Gamma),C_r^*(\Gamma),L(\Gamma)\).
\item[\(p:C_f\to C\)] Full-to-reduced quotient (local symbol in Section~3).
\item[\(w_g,\sigma_g\)] Source group unitaries and lamp-coordinate action.
\item[\(B_g,U_j,W_g\)] \(N\)-valued cocycle, \(r_N(v_j^*)\), \(r_N(B_g^*)\).
\item[\(T_0,\Phi,T,\theta\)] Algebraic compression; full, reduced, normal source maps.
\item[\(P,I,U=P/I\)] Product of matrix blocks, trace ideal, \(C^*\)-quotient.
\item[\(q_\tau,\sigma_C\)] \(P\to U\) quotient and faithful reduced-source copy.
\item[\(H_c,Q_c,j\)] Block Hilbert space, its Calkin algebra, block \(*\)-map \(P\to Q_c\).
\item[\(F_U,F_Q,T_v\)] Trace-GNS return, whole-Calkin return, moment map.
\item[\(\alpha,\Psi\)] Normal corner compression and one fixed tensor state.
\item[\(y,t_0,u_0\)] Same finite coefficients in \(B\odot C,B\odot P,B\odot Q_c\).
\item[\(K_0,h,c=1-h\)] Finite ideal correction and its positive outer cut.
\item[\(u,k\)] Cut selfadjoint tensor and positive max-kernel element.
\end{description}
Local scalar letters and spectral projections are defined where used.
In particular \(D_{\mathrm{dual}}\) below is a duality isomorphism, distinct from \(D=L(\Gamma)\).

The proposed proof has three components. G supplies one corner, one tuple and one normal state from nonnuclearity. M supplies one state-independent moment map, a faithful source copy and exact block partitions from that tuple. D uses these same objects to construct the nonzero positive kernel element over the original \(B\). These components are proved in Sections~2, 3 and~4, respectively; none is invoked merely as an additional hypothesis in the final assembly. Section~5 records their order of use.

\section{Standard inputs and exactness from the hypothesis}
Recall that exactness of \(B\) means that tensoring every \(C^*\) short exact sequence with \(B\) in the minimal norm preserves exactness. Nuclearity means equality of minimal and maximal tensor norms with every \(C^*\)-algebra.
We adopt the following established results at their stated scopes.
\begin{description}
\item[(E)] For every \(C^*\)-algebra \(B\), exactness is equivalent to exactness of
\[
0\longrightarrow B\otimes_{\min}\mathcal K(H_*)
\longrightarrow B\otimes_{\min}\mathcal B(H_*)
\longrightarrow B\otimes_{\min}Q_*\longrightarrow 0.
\]
Kirchberg [1], Theorem~1.1(ii), (iii).
\item[(B)] \(B\) is nuclear if and only if its canonical bidual \(B^{**}\) is injective, also for nonunital and nonseparable \(B\). Choi--Effros [2]. An arbitrary weak closure of a faithful \(C^*\)-representation is not substituted for \(B^{**}\).
\item[(H)] A von Neumann algebra \(F_0\) is finite and injective if and only if for every finite unitary family \(w_1,\ldots,w_k\) and nonzero central projection \(p\),
\[
\Big\|\sum_j p w_j\otimes\overline{p w_j}\Big\|_{\min}=k.
\]
Haagerup [3], Lemma~2.2. We use this only when \(F_0\) is already finite.
\item[(S)] Standard forms exist on arbitrary von Neumann algebras. Every normal positive functional has a unique natural-cone vector, with
\(\|\xi_\phi-\xi_\psi\|^2\le\|\phi-\psi\|\).
The standard implementation \(L(u)r(u^*)\) of an inner automorphism preserves that cone.
Haagerup [4], Theorem~1.6 and Lemma~2.10.
\item[(C)] For a \(\sigma\)-finite \(M\) and faithful normal state \(\chi\), let
\(C_0=M\rtimes_{\sigma^\chi}\mathbb R\), \(\beta\) its dual action, and \(C_2=C_0\rtimes_\beta\mathbb R\).
The algebra \(C_0\) has a normal semifinite faithful trace, and \(C_2\) is normally, faithfully and onto \(*\)-isomorphic to \(M\overline\otimes\mathcal B(L^2(\mathbb R))\), with normal inverse. No factor or separable-Hilbert-space assumption is imposed. [5], [6].
\end{description}

We also use ordinary \(C^*\)/von Neumann theory: Arveson extension into \(\mathcal B(H)\), CP Schwarz, GNS and continuous functional calculus, normal supports, the semifinite/type~III central split, semifinite trace theory and proper infinite halving; faithful min inclusions and maximal commuting representations; source standard left/right commutants and Kaplansky density; Hahn--Banach, Goldstine and weak* compactness. In particular no extension into an arbitrary noninjective commutant is assumed.
Arveson extension here concerns a UCP map from a unital operator system into \(\mathcal B(H)\), with arbitrary \(H\). Kaplansky density concerns the bounded unit ball and positive matrix cones of a \(C^*\)-algebra in its von Neumann closure. Goldstine concerns the weak* density of \(\operatorname{Ball}(B)\) in \(\operatorname{Ball}(B^{**})\) for the specified predual \(B^*\). These applications are spelled out below.

\paragraph{Exactness gateway.}
Maximal tensor products preserve short exact rows for an ideal \(J\) in \(A\).
For the isometry \(B\otimes_{\max}J\to B\otimes_{\max}A\), extend a commuting representation of \(J\) to \(A\) on its essential \(J\)-space by an approximate identity; it still commutes with \(B\). The orthogonal complement contributes zero to \(B\odot J\). A quotient representation annihilating \(B\otimes_{\max}J\) factors through \(A/J\) on the joint essential space. These two universal properties give the maximal exact row; no arbitrary-subalgebra max inclusion is used.

Every canonical max-to-min map is onto: its range is closed and contains the dense algebraic tensor product. Assume \(\chi_B\) is injective. Lift an element \(t\) of the kernel of the minimal \(\mathcal B(H_*)\) quotient row through \(B\otimes_{\max}\mathcal B(H_*)\to B\otimes_{\min}\mathcal B(H_*)\). The maximal quotient image of the lift has zero minimal image, hence is zero by injectivity of \(\chi_B\). Maximal ideal exactness puts the lift in \(B\otimes_{\max}\mathcal K(H_*)\); its minimal image is in \(B\otimes_{\min}\mathcal K(H_*)\). The reverse inclusion is immediate, so the minimal test row is exact. Input~(E) gives exactness of the original \(B\), with no separability or unitization step. This argument does not first assert the hypothesis for a subalgebra or an amplification of \(B\).

\section{General input G in one literal corner}
\paragraph{Proposition G (proposed).}
Let \(V\) be a noninjective von Neumann algebra, with no separability or \(\sigma\)-finiteness assumption. Given \(0<q_0<1\) and \(\varepsilon>0\), there are:
\begin{itemize}
\item a nonzero \(e\in\operatorname{Proj}(V)\), the literal algebra \(N=eVe\);
\item a finite \(F\) with faithful normal tracial state \(\tau\);
\item a normal faithful relative-unital \(*\)-embedding \(i:F\to N\);
\item a relative UCP return \(E:N\to F\), \(Ei=\mathrm{id}_F\), not assumed normal;
\item a finite indexed tuple \(v_1,\ldots,v_n\) in \(i(\mathcal U(F))\);
\item one normal state \(\phi\) on \(N\), not necessarily faithful,
\end{itemize}
such that
\begin{equation}\tag{G}
q=\left\|n^{-1}\sum_j v_j\otimes\overline{v_j}\right\|_{N\otimes_{\min}\overline N}\le q_0,\quad
\max_j\|\phi\circ\operatorname{Ad}(v_j^*)-\phi\|_{N_*}<\varepsilon.
\end{equation}
Here \(\overline N\) is the conjugate \(C^*\)-algebra; \(q\) is an operator norm, not its square. With the standard form \(L_N,J_N\) and \(r_N(x)=J_N L_N(x^*)J_N\),
\begin{equation}\tag{G-CORR}
r=n^{-1}\sum_j\operatorname{Re}\langle L_N(v_j)r_N(v_j^*)\xi_\phi,\xi_\phi\rangle>1-\varepsilon/2.
\end{equation}
Thus \(q_0=1/4\) and \(\varepsilon=1/32\) give \(q<1/3\) and \(r>63/64\) on the same data.

\paragraph{Proof.} Adopt the standard inputs (B), (H), (S), (C) from Section~1.

\subsection{Localization without a separable representation}
Represent \(V\) faithfully normally as \(Q\subset\mathcal B(K)\), allowing arbitrary \(K\).
Injectivity is equivalent to a UCP map \(\mathcal B(K)\to Q\) fixing \(Q\).
The forward implication extends \(\mathrm{id}_Q\); the converse extends a map into \(\mathcal B(K)\) and then composes. This also proves independence of the faithful normal copy.
The point-ultraweak compact space \(\operatorname{UCP}(\mathcal B(K),Q)\) is nonempty: it contains \(X\mapsto\varpi(X)1\). Linearity, unit and all matrix-positive constraints are closed in the product of weak* compact coordinate balls. Noninjectivity makes the closed constraints \(\nu(T(x))=\nu(x)\), for all selfadjoint contractions \(x\in Q\) and normal states \(\nu\), have empty intersection.
A finite subfamily already has empty intersection. Let \(s\) be the support of the sum of those finitely many states. Their normalized sum is a faithful normal state on \(sQs\). If \(sQs\) were injective, a return \(R_s\) on \(\mathcal B(sK)\) would make
\[
T(X)=R_s(sXs)+\varpi(X)(1-s)
\]
a UCP map satisfying all selected tests, since \(\nu(sxs)=\nu(x)\). Contradiction. Pull back \(s\) to \(V\). This is a literal \(\sigma\)-finite noninjective corner, without selecting factor fibers or a common exceptional null set. Split its semifinite and type~III central parts. One is noninjective: returns for both parts would give a return for their direct sum. Call that part \(M=p_M Vp_M\).

Another compactness fact will be used twice. If \(p_a\) increases strongly to \(1\) in \(Q\) and every \(p_aQp_a\) is injective, the maps
\[
 T_a(X)=R_a(p_a Xp_a)+\varpi(X)(1-p_a)
\]
have a cofinal point-ultraweak cluster map fixing \(Q\), since on \(x\in Q\) their values are \(p_a xp_a+\varpi(x)(1-p_a)\to x\). Thus \(Q\) is injective.
For an n.s.f.\ trace, finite-trace projections are directed under joins
(\(\operatorname{Tr}(p\vee q)\le\operatorname{Tr}(p)+\operatorname{Tr}(q)\)) and have supremum \(1\) by semifiniteness. Therefore a noninjective semifinite algebra has a noninjective finite-trace corner. No sequence is substituted for this general net.

If \(M\) is semifinite, choose \(f\) with \(0<\operatorname{Tr}(f)<\infty\) and \(fMf\) noninjective. Take \(F_0=N_0=fVf\), \(e_0=f\), \(i_0=R_0=\mathrm{id}\) and \(\tau_0=\operatorname{Tr}|_{F_0}/\operatorname{Tr}(f)\). Continue with Section~2.4.

\subsection{Type III: exact copies, bounded returns and a literal target}
Fix a faithful normal \(\chi\) of \(M\) for its modular action only. Put \(j_0:M\to C_0\), \(j_1:C_0\to C_2\), the normal faithful canonical unital embeddings; \(\beta\) is the first dual action, \(\gamma\) the second dual action.
Their fixed algebras are \(j_0(M)\) and \(j_1(C_0)\), respectively. For clarity, this fixed-copy identification needs no separable Hilbert space: in the regular gauge of a spatially implemented \(\mathbb R\)-action \(\alpha\) on \(Q\), the product is
\[
\{Q\otimes1,V_t\otimes\lambda_t\}'',\qquad \lambda_t f(s)=f(s-t).
\]
The dual is \(\operatorname{Ad}(1\otimes m_h)\), \(m_h f(s)=\exp(ihs)f(s)\). Every product element commutes with scalar translations. A dual-fixed element also commutes with scalar modulations, hence equals \(x\otimes1\). For \(y\in Q'\), the multiplier \(Y_y(s)=V_s y V_s^*\) commutes with both generating families. Thus \([x,V_s y V_s^*]\) has zero coefficients almost everywhere for each fixed vector pair. Those coefficients are continuous; evaluate at \(0\) to obtain \([x,y]=0\). Hence \(x\in Q\).
Standard form gives the required strongly continuous spatial implementation for each action.

For \(\eta=\beta\) or \(\gamma\), define the normal UCP averages
\[
 A_L^\eta(X)=(2L)^{-1}\int_{-L}^L\eta_t(X)\,dt.
\]
The ultraweak integral is normal: predual orbits are norm-continuous, as follows from strong unitary implementation and trace-class approximation in a normal representation. Moreover
\[
 \|\eta_h(A_L^\eta(X))-A_L^\eta(X)\|\le |h|\,\|X\|/L.
\]
A cofinal point-ultraweak cluster limit \(P_\eta\) is UCP into the fixed copy and fixes that entire copy. Its normality is not inferred. Set
\begin{equation}\tag{RET}
\begin{aligned}
 E_0&=j_0^{-1}P_\beta:C_0\to M, & E_0j_0&=\mathrm{id}_M,\\
 E_1&=j_1^{-1}P_\gamma:C_2\to C_0, & E_1j_1&=\mathrm{id}_{C_0}.
\end{aligned}
\end{equation}
For every \(a,b\in C_0\) and \(Y\in C_2\), the multiplicative domain gives
\begin{equation}\tag{MOD}
 E_1(j_1(a)Yj_1(b))=a E_1(Y)b.
\end{equation}
The map \(E_0\) is used here only to reflect noninjectivity: if \(C_0\) were injective, extension into \(C_0\) followed by \(E_0\) would make \(M\) injective. Thus \(C_0\) is noninjective. This is a retract argument, not inheritance by arbitrary subalgebras.

Let \(D_{\mathrm{dual}}:C_2\to M\overline\otimes\mathcal B(L^2(\mathbb R))\) be the normal onto duality isomorphism. In \(\sigma\)-finite type~III \(M\), halve its unit into two projections equivalent to it, retain the half of smaller \(\chi\)-mass, and repeat. Remainder masses tend to \(0\); faithfulness and normality of \(\chi\) make their meet zero. The removed pieces are orthogonal and equivalent to \(1_M\), with sum \(1_M\).
Choose isometries \(s_l\) with these ranges. The formula \(U(\eta\otimes\delta_l)=s_l\eta\) defines a unitary \(K_M\otimes\ell^2\to K_M\). For \(A_M(X)=UXU^*\),
\[
 A_M(x\otimes e_{ij})=s_i x s_j^*\in M.
\]
Finite matrix compressions and strong closure give its range in \(M\). Conversely \(U^*xU\) has entries \(s_i^*xs_j\in M\), and the same compression argument puts it in \(M\overline\otimes\mathcal B(\ell^2)\). Thus \(A_M\) is normal and onto, with normal inverse. A fixed unitary \(L^2(\mathbb R)\to\ell^2\) and \(D_{\mathrm{dual}}\) now give a normal onto \(*\)-isomorphism \(\Xi:C_2\to M\).
It need not fix \(j_1j_0(M)\); it does map \(1_{C_2}\) to \(1_M=p_M\).

\subsection{A common finite corner and the second return}
The following refinement also distinguishes the two returns on the same finite corner. It is available within every selected type~III core.
Write its implementing group \(\ell_t=\exp(itA)\), with \(\beta_s(\ell_t)=\exp(ist)\ell_t\), and put \(p_n=1_{[-n,n]}(A)\). The \(p_n\) increase strongly to \(1\) but need not be finite. Spectral calculus and the length of an interval intersection give
\begin{equation}\tag{BAND}
\begin{aligned}
A_L^\beta(p_n)&=g_{n,L}(A),\\
g_{n,L}(a)&=\frac{\operatorname{length}([-L,L]\cap[-n-a,n-a])}{2L},\\
0\le A_L^\beta(p_n)&\le \min(1,n/L)1;\quad \text{hence } E_0(p_n)=0.
\end{aligned}
\end{equation}
Section~2.1 implies some \(p_N C_0p_N\) is noninjective. The restricted core trace is n.s.f.\ there, so its finite-trace exhaustion supplies \(f\le p_N\), \(0<\operatorname{Tr}(f)<\infty\), with \(F_0=fC_0f\) noninjective.
Since \(p_n\) increases to \(1\) but \(E_0(p_n)=0\), \(E_0\) is nonnormal. Positivity also gives \(E_0|_{F_0}=0\). Set \(\tau_0=\operatorname{Tr}|_{F_0}/\operatorname{Tr}(f)\) and
\begin{equation}\tag{CORNER}
\begin{aligned}
e_0 &=\Xi(j_1(f)), & N_0&=e_0Ve_0=e_0Me_0,\\
i_0 &: F_0\to N_0, & i_0(a)&=\Xi(j_1(a)),\\
R_0 &: N_0\to F_0, & R_0(X)&=E_1(\Xi^{-1}(X)).
\end{aligned}
\end{equation}
Onto-ness, not merely faithfulness, proves \(\Xi^{-1}(N_0)=j_1(f)C_2j_1(f)\).
Thus (MOD) proves \(R_0(X)\in F_0\), \(R_0(e_0)=f\), \(R_0i_0=\mathrm{id}_{F_0}\).
The map \(i_0\) is normal faithful, sends the unit \(f\) to \(e_0\), and sends each projection of \(F_0\) to a projection of \(N_0\). The map \(R_0\) is UCP relative to these units; no normality is claimed. The projection \(f\) is finite in \(C_0\), whereas its nonzero image \(e_0\) is a type~III projection in \(M\). Finiteness is relative to the ambient algebra and was not required to pass to \(M\).

The comparison uses exactly this \(f\) and these maps:
\begin{equation}\tag{TEST}
\begin{gathered}
 E_0(f)=0,\qquad E_1(j_1(f))=f,\\
 \mathrm{Bad}_0=E_0E_1\Xi^{-1}:N_0\to M \text{ is identically zero},\\
 R_0=E_1\Xi^{-1}:N_0\to F_0\text{ is relative unital}.
\end{gathered}
\end{equation}
The shorter \(E_0\Xi^{-1}\) is ill-typed (\(C_2\) versus \(C_0\)). The type-correct \(\mathrm{Bad}_0\) fails already at \(\mathrm{Bad}_0(e_0)=0\) and cannot be normalized to a state. There is no incompatibility: \(\beta\) moves the core elements while \(\gamma\) fixes all of \(j_1(C_0)\). The maps \(E_0\) and \(E_1\) have different domains and fixed copies.

\subsection{One central cut and internal powers}
Apply input~(H) to finite noninjective \(F_0\). There are \(w_j\in\mathcal U(F_0)\) and a nonzero \(p\in\operatorname{Proj}(Z(F_0))\) with \(u_j=pw_j\) in \(F=pF_0p\) and
\[
a=\left\|k^{-1}\sum_j u_j\otimes\overline{u_j}\right\|_{F\otimes_{\min}\overline F}<1.
\]
Faithful corner inclusions preserve the min norm used in input~(H). Define
\[
e=i_0(p),\quad N=eVe=eN_0e,\quad i=i_0|_F,\quad E=R_0|_N,\quad \tau=\tau_0|_F/\tau_0(p).
\]
The projection \(p\) need not be central in \(C_0\); \(e\) need not be central in \(N_0\) or \(V\). Nevertheless the fixed-subalgebra module rule gives
\begin{equation}\tag{G-CUT}
\begin{gathered}
 E:N\to F,\quad E(e)=p,\quad Ei=\mathrm{id}_F,\\
 E(i(a)X i(b))=aE(X)b\quad(a,b\in F,\ X\in N).
\end{gathered}
\end{equation}
In the semifinite case these formulas hold with the identity return. In the type~III case \(E_0(p)=0\) but \(E(e)=p\); (TEST) survives this same cut.
If \(\psi_0=\tau_0 R_0\), then \(\psi=\tau E=\psi_0|_N/\psi_0(e)\), with \(\psi_0(e)=\tau_0(p)>0\).
This is the one normalized restriction. After it, \(e,N,F,i,E,\tau,\psi\) remain fixed. For every \(w\in\mathcal U(F)\), traciality and (G-CUT) give
\begin{equation}\tag{INV}
\psi(i(w)^*X i(w))=\psi(X)\quad(X\in N).
\end{equation}

Put \(z_j=i(u_j)\), \(d=k^{-1}\sum_j z_j\otimes\overline{z_j}\). Normal faithfulness of \(i\) implies faithful \(C^*\)-inclusion, so \(\|d\|_{\min}=a\). Choose \(m\ge1\) with \(a^{2m}\le q_0\). Keep all indexed words, including repeated values,
\[
v_J=z_{j_1}^*z_{j_2}\cdots z_{j_{2m-1}}^*z_{j_{2m}},\qquad n=k^{2m}.
\]
Their paired average is exactly \((d^*d)^m\), whose norm is \(a^{2m}\), because \(d^*d\) is positive. No tensor power of \(N\), matrix amplification or new corner is introduced. Every word lies in \(i(F)\), so (INV) holds for it. The return is used multiplicatively only on its fixed embedded algebra. The whole final list and its length are now fixed, before selecting \(\phi\).

\subsection{One normal state from finite defects}
Let \(\alpha_j(X)=v_j^*Xv_j\). These are normal automorphisms of \(N\). The set
\[
C=\{(\phi\circ\alpha_j-\phi)_j:\phi\text{ a normal state of }N\}
\]
is convex in \((N_*^{\mathrm h})^n\), with maximum norm and real dual \((N^{\mathrm{sa}})^n\).
If \(0\) is outside its norm closure, strict separation supplies \(H_j=H_j^*\) and \(b>0\) with \(\phi(H)\ge b\) for every normal state \(\phi\), where
\[
H=\sum_j(v_j^*H_jv_j-H_j).
\]
Normal states detect order, hence \(H\ge be\). But the state \(\psi\) in the full dual \(N^*\) has \(\psi(H)=0\) by (INV), contradicting \(\psi(H)\ge b\). Consequently one normal \(\phi\) makes every defect \(<\varepsilon\). This needs invariance on all \(N\), not just on \(i(F)\). It assumes neither normality of \(\psi\) nor norm approximation of \(\psi\) by normal states. It asserts neither exact invariance of \(\phi\) nor existence of a normal invariant limiting state.

For completeness this distinction is forced in the actual type~III construction. With second implementers \(\exp(it\widehat A)\), put \(Q_l=1_{[-l,l]}(\widehat A)\) in \(C_2\). The same band calculation gives \(E_1(Q_l)=0\).
The positive contractions \(X_l=\Xi(j_1(p)Q_l j_1(p))\) increase strongly to \(e\); they need not be projections. Equation (MOD) gives \(E(X_l)=0\). Thus \(\psi(X_l)=0\) but every normal \(\phi\) has \(\phi(X_l)\to1\). Testing \(2X_l-e\) gives \(\|\phi-\psi\|_{N^*}=2\). In particular \(E\) is not normal in this branch.
Finite predual defect approximation and this distance-\(2\) assertion have different quantifiers and are consistent. No cone vector is used for \(\psi\).

For the selected normal \(\phi\), the operator \(W_j=L_N(v_j)r_N(v_j^*)\) has \(W_j^*L_N(X)W_j=L_N(v_j^*Xv_j)\), and preserves the natural cone. Input~(S) gives
\[
2\bigl(1-\operatorname{Re}\langle W_j\xi_\phi,\xi_\phi\rangle\bigr)
=\|W_j\xi_\phi-\xi_\phi\|^2
\le\|\phi\circ\alpha_j-\phi\|<\varepsilon,
\]
which proves the correlation assertion (G-CORR) of Proposition~G.
Nonfaithfulness triggers no support cut, state mixture, or new modular core.
The states \(\chi,\tau,\psi,\phi\) have distinct roles throughout.

\subsection{The canonical bidual and nonunital \texorpdfstring{\(B\)}{B}}
For nonnuclear \(B\), use input~(B) with \(V=B^{**}\). Every projection chosen above is in \(V\); \(\Xi\) absorbs the type~III amplification before the final \(e\) and \(\phi\) are selected. Normal compression \(c_e:V\to eVe\), \(c_e(X)=eXe\), sends \(1_V\) to \(e\) and fixes each projection below \(e\), but is not generally multiplicative. The functional \(\phi_V=\phi c_e\) is normal on \(V\). Its restriction to a nonunital \(B\) has norm \(1\): the positive contractive approximate-identity net converges strongly to \(1_V\). This does not claim \(v_j\in B\) or perform the later transfer D.

\paragraph{Quantifier order.}
Starting with \(V\), choose one \(\sigma\)-finite noninjective \(M\); then, if needed, the core, copies, returns and onto \(\Xi\); then finite-trace noninjective \(F_0\) and \(N_0\); then one central \(p\) and fixed \(F,N,i,E,\tau,\psi\); then one full word list; and finally one \(\phi\).
The two returns can be fixed before \(f\); (BAND) holds for every cluster choice of \(E_0\), and \(E_1\) fixes every core element before any \(f\) is selected. No normality of a return, centrality in the ambient target, separable representation, or state on a leftover amplification is an input.

This proves Proposition~G.


\section{The cyclic-lamp moment map M}

\paragraph{Proposition M (proposed, state-independent moment bridge).}
Let \(N\) be any von Neumann algebra in standard form and let \(v_1,\ldots,v_n\) be relative unitaries with
\[
 q=\left\|n^{-1}\sum_j v_j\otimes\overline{v_j}\right\|_{N\otimes_{\min}\overline N}<1/3,\qquad n\ge1.
\]
Put \(G=F(s,t)\), \(\Gamma=(\mathbb Z/(n^2))^{(G)}\rtimes G\), with finite-support lamps, \(C=C_r^*(\Gamma)\), and \(D=L(\Gamma)\). Identify each cyclic lamp's \(n^2\) spectral atoms with pairs in \(\{0,\ldots,n-1\}^2\); partition labels \(j\) run from \(1\) to \(n\). Its row and column partitions at \(h\) are \(p_{h,s,j},p_{h,t,j}\); let \(d_{u,j}=w_u p_{1,u,j}\), \(u=s,t\).
Fix any one countable sequence of group homomorphisms \(\rho_l:\Gamma\to U(d_l)\) whose normalized traces tend to the regular \(\Gamma\) trace, and one free ultrafilter \(\omega\). Set
\[
 \begin{gathered}
 P=\prod_l M_{d_l},\quad I=\{b:\lim_\omega\operatorname{tr}_l(b_l^*b_l)=0\},\quad U=P/I,\\
 H_c=\bigoplus_l\mathbb C^{d_l},\quad Q_c=B(H_c)/K(H_c),\quad q_\tau:P\to U\ \text{the quotient}.
 \end{gathered}
\]
Here \(j:P\to Q_c\) is block-diagonal inclusion followed by the Calkin quotient.
Let \(S_u\) and \(P_{u,j}\) be the model shifts and lamp partitions and put \(b_{u,j}=S_u P_{u,j}\). The shifts are unitaries and the \(P_{u,j}\) form exact orthogonal partitions in \(P\). There is a faithful unital embedding
\[
 \sigma_C:C\to U,\qquad \sigma_C(d_{u,j})=q_\tau(b_{u,j}),
\]
and there is one UCP map \(T_v:Q_c\to L_N(N)'\) such that, simultaneously,
\begin{equation}\tag{M}
 T_v(j(b_{u,j}))=r_N(v_j^*)/n,\qquad T_v(j(i))=0\quad\text{for every }i\in I.
\end{equation}
No target state, original \(B\), or later ideal cut is involved in choosing this map. It is not asserted multiplicative, normally lifted to \(B(H_c)\), or onto the target commutant. Its specified values, not just range containment, are essential. The proof below supplies such a sequence and proves the conclusion for any one fixed sequence satisfying the stated trace condition.

\noindent\textit{Proof.}

\subsection{The source group and its cocycle}

The lamp group consists of finite-support functions \(a:G\to\mathbb Z/(n^2)\). With \((g\cdot b)(h)=b(g^{-1}h)\), its semidirect product law and inverse are
\[
 (a,g)(b,h)=(a+g\cdot b,gh),\qquad (a,g)^{-1}=(-g^{-1}\cdot a,g^{-1}).
\]
The Fourier dual of the lamp group is \(\Omega=(\{0,\ldots,n-1\}^2)^G\) with uniform product measure. Put \(R=L^\infty(\Omega)\). At site \(h\), \(p_{h,s,j}\) and \(p_{h,t,j}\) are the indicators of row \(j-1\) and column \(j-1\). The cyclic character has value \(\exp(2\pi i(na+b)/n^2)\) on the pair \((a,b)\). Thus the lamp is the original cyclic group of order \(n^2\), not the group \(C_n\times C_n\).
Write \(\sigma_g(p_{h,u,j})=p_{gh,u,j}\), \(w_g a w_g^*=\sigma_g(a)\). Then \(D=R\rtimes_{\mathrm{reg}} G=L(\Gamma)\), \(A=\mathbb C[\Gamma]\) consists of finite sums \(x=\sum_g a_g w_g\) with finite cylinder \(a_g\), and \(C=C_r^*(\Gamma)\). Here
\((a w_g)(b w_h)=a\sigma_g(b)w_{gh}\) and \((a w_g)^*=\sigma_{g^{-1}}(a^*)w_{g^{-1}}\).
Its faithful regular trace satisfies \(\tau_D(a w_g)=0\) for \(g\ne1\) and \(\tau_D(a)=\int a\,d\mu\); thus \(\|x\|_2^2=\sum_g\|a_g\|_2^2\). Finite cyclic characters form the exact regular group orthonormal basis, not merely an approximate family.

For reduced \(g=u_1^{\varepsilon_1}\cdots u_l^{\varepsilon_l}\), let \(g_i\) be its prefixes, \(g_0=1\). Edge \(i\) has key \((g_i,u_i)\) if \(\varepsilon_i=+1\) and key \((g_{i-1},u_i)\) if \(\varepsilon_i=-1\). The key \((h,u)\) denotes the unoriented edge \(hu^{-1}\to h\) in either direction. Put
\[
 Q_{g,J}=\prod_i p_{\mathrm{key}_i,j_i},\quad
 V_{g,J}=v_{j_1}^{\varepsilon_1}\cdots v_{j_l}^{\varepsilon_l},\quad
 B_g=\sum_J Q_{g,J}\otimes V_{g,J}.
\]
These are finite cylinder functions with unitary values in \(N\). Reversing an edge uses the same partition, not a fresh label. On an adjacent backtrack the two projections force equal labels, and adjacent \(vv^*\) or \(v^*v\) cancels without moving other factors. Translation sends \((h,u)\) to \((gh,u)\). Repeated adjacent cancellation therefore proves for all words
\begin{equation}\tag{COC}
 B_{gh}=B_g\sigma_g(B_h),\qquad B_{g^{-1}h}=\sigma_{g^{-1}}(B_g^*B_h).
\end{equation}
This includes concatenations with an arbitrarily long cancellation.

On \(L^2(R)\) let \(L_R(a)\) be left multiplication and \(r_R(a)\) right multiplication; they coincide because \(R\) is abelian. The unitary \(S_g\) is specified by \(S_g\widehat a=\widehat{\sigma_g(a)}\). On \(H_{\mathrm{raw}}=L^2(R)\otimes H_N\), put \(\mathcal R=r_R\otimes r_N\), a linear anti-representation on the finite \(R\)-valued tensor sums used here, and define
\[
 \Pi_g=\mathcal R(B_g^*)(S_g\otimes 1),\qquad J_0\eta=1\otimes\eta.
\]
Anti-multiplication and (COC) give \(\Pi_g\Pi_h=\Pi_{gh}\). Covariance with \(L_R\) then gives a group \(*\)-representation \(\pi_{\mathrm{alg}}(a w_g)=L_R(a)\Pi_g\) and hence \(\pi_f\) of the full \(C_f=C^*(\Gamma)\). Its UCP compression \(\Phi=J_0^*\pi_f(\cdot)J_0\) has range in \(L_N(N)'\) and on \(A\) is
\begin{equation}\tag{FWD}
 T_0(a w_g)=\int a(\omega)r_N(B_g(\omega)^*)\,d\mu
 =\sum_J\tau_R(aQ_{g,J})r_N(V_{g,J}^*).
\end{equation}
The integral is a finite sum; no measurability of nonseparable operator spaces is assumed. For \(d_{u,j}=w_u p_{1,u,j}=p_{u,u,j}w_u\) it gives
\begin{equation}\tag{MOM}
 T_0(d_{u,j})=r_N(v_j^*)/n.
\end{equation}

\subsection{Positivity at every matrix degree}

If \(x=a w_g,y=b w_h\), then \(x^*y=\sigma_{g^{-1}}(a^*b)w_{g^{-1}h}\). Using (COC) and invariance of \(\mu\),
\begin{equation}\tag{GRAM}
 T_0(x^*y)=\int\overline a b\,r_N(B_h^*B_g)\,d\mu
 =\int A_x^*A_y\,d\mu,\qquad A_x=a r_N(B_g^*).
\end{equation}
The matrix-entry order is \(B_h^*B_g\) because \(r_N\) is anti-multiplicative. For arbitrary rectangular \(Z=[z_{li}]\) with \(z_{li}=\sum_g a_{li,g} w_g\) put \(F_{li}=\sum_g a_{li,g}r_N(B_g^*)\). Then for every finite \(r,k\),
\begin{equation}\tag{CP-A}
 \begin{aligned}
 T_0^{(k)}(Z^*Z)&=\int F^*F\,d\mu,\\
 \langle T_0^{(k)}(Z^*Z)\eta,\eta\rangle
 &=\sum_l\int\left\|\sum_i F_{li}\eta_i\right\|^2d\mu\ge0.
 \end{aligned}
\end{equation}
Choose one finite set \(S\) of sites containing all coefficients' lamp supports and all words' edge keys. The integral is precisely
\[
 n^{-2|S|}\sum_{\omega_S} F(\omega_S)^*F(\omega_S).
\]
Each site receives one pair of labels across every row, column and word. Shared edges share labels; refining \(S\) just sums unused probabilities to \(1\). Thus all matrix sizes concern the same map, with every cross term kept.
This proves algebraic matrix-square positivity and is compatible with the full compression \(\Phi\). It does not by itself prove reduced positivity or normality; those are the next two steps.

\subsection{A uniform Hilbert--Schmidt frame estimate}

\paragraph{Lemma HS (proposed Hilbert--Schmidt frame estimate).}
Let \(H\) be a nonzero complex Hilbert space, with no separability hypothesis, and \(U_1,\ldots,U_n\) unitaries, \(n\ge1\), with repetitions allowed. Inner products are linear in the first variable; \(\lvert\eta\rangle\langle\zeta\rvert\) sends \(x\) to \(\langle x,\zeta\rangle\eta\). On \(S_2(H)\), the ambient Hilbert--Schmidt space, define
\[
 P_+(X)=n^{-1}\sum_j U_jXU_j^*,\qquad P_-(X)=n^{-1}\sum_j U_j^*XU_j .
\]
Assume \(q=\|P_+\|_{S_2\to S_2}<1/3\) and set \(C_q^2=(1+q)/(1-3q)\).
\(P_-=P_+^*\) on \(S_2\), so its norm is also \(q\). These are averages of unitary conjugations on \(S_2\); each has norm \(1\) on \(B(H)\). No state or trace on a target von Neumann algebra is assumed. The trace defining \(S_2\) is the ambient one.

Let \(G=F(s,t)\), \(\Omega=(\{0,\ldots,n-1\}^2)^G\) with uniform product probability. At a site \(h\) let \((R_h,C_h)\) be its coordinate pair. This is the spectrum of one cyclic lamp of order \(n^2\), labeled by
\[
 \zeta_h=\exp(2\pi i(nR_h+C_h)/n^2).
\]
The characters \(\chi_\alpha=\prod_h\zeta_h^{\alpha(h)}\), for finite-support \(\alpha:G\to\mathbb Z/(n^2)\), are its original cyclic Fourier basis. Relabeling the uniform atom by a pair does not change the underlying cyclic lamp group.

For a reduced \(g=u_1^{\varepsilon_1}\cdots u_l^{\varepsilon_l}\), let \(g_i\) be its successive prefixes, \(g_0=e\), \(u_i\in\{s,t\}\), \(\varepsilon_i\in\{+1,-1\}\). Its \(i\)-th slot is \((h_i,u_i)\), where
\[
 h_i=g_i\ \text{if }\varepsilon_i=+1,\qquad
 h_i=g_{i-1}\ \text{if }\varepsilon_i=-1.
\]
The key \((h,u)\) denotes the unoriented edge \(hu^{-1}\)--\(h\), with that fixed positive orientation. Use \(j_i=1+R_{h_i}\) for \(u_i=s\) and \(j_i=1+C_{h_i}\) for \(u_i=t\), and define
\[
 W_g=U_{j_1}^{\varepsilon_1}\cdots U_{j_l}^{\varepsilon_l},\qquad W_e=I.
\]
All words use the same \(\Omega\). Reversing an edge reads the same slot and uses the adjoint. Common edges and prefixes across words share labels.

Let \(x=\sum_g a_gw_g\) be a finite sum with arbitrary complex cylinder functions \(a_g\). These are precisely the finite group-algebra sums for \(\Gamma=C_{n^2}^{(G)}\rtimes G\). For its regular trace,
\[
 \|x\|_2^2=\sum_g\|a_g\|_{L^2(\Omega)}^2.
\]
Define \(T_0(x)=\sum_g\mathbb E[a_g W_g]\), a finite sum of bounded operators. Then for every such \(x\) and every \(\eta\in H\),
\begin{equation}\tag{HS1}
 \|T_0(x)\eta\|^2\le C_q^2\|x\|_2^2\|\eta\|^2,
\end{equation}
and therefore
\begin{equation}\tag{HS2}
 \|T_0(x)\|_{\mathrm{op}}\le C_q\|x\|_2.
\end{equation}
The same constant works for all \(\eta\), including choices depending on \(x\).
\(T_0\) has one bounded linear extension \(L^2(L(\Gamma))\to B(H)\). This assertion places no order or completely positive map structure on the Hilbert space \(L^2\).

\noindent\textit{Proof.}

\subsubsection{The joint law on one reduced path}

A reduced path in the free Cayley tree is simple: equality of two prefix vertices would make a nonempty reduced subword equal to the identity. An edge-head \(h_i\) is one endpoint of the \(i\)-th edge. Two heads can coincide only at a common vertex of consecutive edges. Such a coincidence occurs exactly for a positive then negative step. Their generators must differ, otherwise they cancel. Thus the repeated site contributes its row once and column once. Positive-positive and negative-negative consecutive steps have distinct heads; negative-positive cannot repeat a head without a forbidden cancellation. Nonadjacent heads cannot coincide.

If the \(l\) slots occupy \(d\) sites, prescribing all \(l\) labels leaves exactly \(n^{2d-l}\) atoms among \(n^{2d}\) at those sites. Every site is uniform and different sites are independent by definition of \(\Omega\). The \(l\) labels of this one path consequently have the full product-uniform law. This is joint independence along a path, not only pairwise independence. It does not assert independence of labels for two different paths. The reversed inverse word is also reduced, so the same argument covers it, in its own actual slot order; no independent resampling of its labels is made.

\subsubsection{Finite cyclic Fourier identity}

Fix \(\eta\) and \(g\). Let \(Z_g\) be the finite set of its used sites, \(d=|Z_g|\), \(N_g=n^{2d}\), and \(v_a=W_g(a)\eta\) for \(a\) in the finite atom set \(\Omega_{Z_g}\). For each of its \(N_g\) original cyclic characters \(\chi\) define
\[
 z_{g,\chi}=T_0(\chi w_g)\eta=N_g^{-1}\sum_a\chi(a)v_a .
\]
A character with a nonzero coordinate outside \(Z_g\) has zero image, because that coordinate is independent and its character has mean zero. For the characters on \(Z_g\), exact finite Fourier orthogonality gives
\[
 \sum_\chi \chi(a)\overline{\chi(b)}=N_g\delta_{a,b}.
\]
Expanding each rank-one operator, with all sums finite, proves
\begin{equation}\tag{HS3}
 \begin{aligned}
 S_{\eta,g}&:=\sum_\chi\lvert z_{g,\chi}\rangle\langle z_{g,\chi}\rvert
 =N_g^{-1}\sum_a\lvert v_a\rangle\langle v_a\rvert\\
 &=\mathbb E[W_gR_\eta W_g^*]
 =P_{\varepsilon_1}\cdots P_{\varepsilon_l}(R_\eta),\qquad
 R_\eta=\lvert\eta\rangle\langle\eta\rvert .
 \end{aligned}
\end{equation}
The last equality follows by integrating the innermost label first, using Section~3.3.1's full joint law. It is ordered composition; no commutation or symmetry of \(P_+\) and \(P_-\) is assumed. In particular reverse signs use \(P_-\), not \(P_+\), and reversing the word reverses the appropriate order. For the empty word, only the trivial lamp character survives and \(S_{\eta,e}=R_\eta\).

\subsubsection{Summability and the positive limit}

Since \(\|R_\eta\|_{S_2}=\|\eta\|^2\) and \(\|P_+\|_{S_2}=\|P_-\|_{S_2}=q\), equation (HS3) yields
\begin{equation}\tag{HS4}
 \|S_{\eta,g}\|_{S_2}\le q^{|g|}\|\eta\|^2 .
\end{equation}
Here \(q\) is an operator norm of a channel on \(S_2\), not its squared norm. The frame already contains squared vector coefficients. Consequently its norm is bounded by \(q^l\), not by a gratuitous replacement \(q^{2l}\).
There are \(4\) choices for the first letter and \(3\) for each later reduced letter, so the number of length-\(l\) words is \(4\cdot3^{l-1}\) for \(l\ge1\). Hence
\begin{equation}\tag{HS5}
 \sum_g\|S_{\eta,g}\|_{S_2}
 \le\left[1+\sum_{l\ge1}4\cdot3^{l-1}q^l\right]\|\eta\|^2
 =\left[1+\frac{4q}{1-3q}\right]\|\eta\|^2=C_q^2\|\eta\|^2 .
\end{equation}
This is absolute convergence of the word-grouped series in the Banach space \(S_2(H)\). Its sum \(S_\eta\) is positive, since finite positive sums converge in Hilbert--Schmidt and hence operator norm. Every finite word sum is \(\le S_\eta\), and
\begin{equation}\tag{HS6}
 \|S_\eta\|_{\mathrm{op}}\le\|S_\eta\|_{S_2}\le C_q^2\|\eta\|^2 .
\end{equation}
The length-tail bound is explicit:
\begin{equation}\tag{HS7}
 \left\|\sum_{|g|>L} S_{\eta,g}\right\|_{\mathrm{op}}
 \le\frac{4q(3q)^L}{1-3q}\|\eta\|^2,\quad L\text{ a nonnegative integer}.
\end{equation}
There is no exchange of an infinite Bochner integral with a word sum: all expectations in (HS3) were finite before (HS5) was used.

Absolute summability in (HS5) does not imply absolute summability of the ungrouped rank-one terms. For \(\eta\) nonzero, unitarity and (HS3) give at each word
\begin{equation}\tag{HS8}
 \sum_\chi\bigl\|\lvert z_{g,\chi}\rangle\langle z_{g,\chi}\rvert\bigr\|_{S_2}
 =\sum_\chi\|z_{g,\chi}\|^2
 =\operatorname{Tr}(S_{\eta,g})=\mathbb E\|W_g\eta\|^2=\|\eta\|^2.
\end{equation}
Its sum over \(g\) diverges, even when \(q<1/3\). Likewise \(\mathbb E\|W_g\eta\|^2\) is not bounded by \(q^{|g|}\|\eta\|^2\) for a nonempty word. Pointwise, \(W_g\) is unitary:
\begin{equation}\tag{HS9}
 \|W_g(\omega)\eta\|=\|\eta\|,\qquad
 \|W_g(\omega)R_\eta W_g(\omega)^*\|_{S_2}=\|\eta\|^2 .
\end{equation}
Only the norm of the averaged positive frame in (HS4) decays. These identities are exact counterexamples to the two false strengthenings, not to (HS1).
Infinite trace of \(S_\eta\) is compatible with its finite Hilbert--Schmidt norm.

For clarity, the ungrouped positive rank-one series does converge as a net of finite subsets in operator norm, without absolute Hilbert--Schmidt convergence.
After a finite subset includes all active characters of all \(|g|\le L\), its missing positive terms are dominated by the full tail in (HS7). Any extra retained characters only improve this bound. This uses the finiteness of each active Fourier block, positivity and (HS7), not an arbitrary exchange of an absolutely convergent double series. Scalar nonnegative sums can also be regrouped by Tonelli; they will be bounded explicitly below.

\subsubsection{All complex coefficients and mixed Gram terms}

For a finite subset \(F\) of pairs \((g,\chi)\), put \(z_\alpha=z_{g,\chi}\). Every partial Fourier block is a positive sub-sum of (HS3). Therefore
\begin{equation}\tag{HS10}
 S_{\eta,F}=\sum_{\alpha\in F}\lvert z_\alpha\rangle\langle z_\alpha\rvert\le S_\eta.
\end{equation}
Let \(A_{\eta,F}:\ell^2(F)\to H\) be the complex-linear column operator with these columns. Then \(A_{\eta,F}A_{\eta,F}^*=S_{\eta,F}\), whereas \(A^*A\) is the full mixed Gram matrix. No cross term for shared prefixes, inverse words, common lamp coordinates, or equal target outputs has been set to zero.
For every complex coefficient vector \(c\),
\begin{equation}\tag{HS11}
 \left\|\sum_{\alpha\in F} c_\alpha z_\alpha\right\|^2
 \le\|A_{\eta,F}\|^2\sum_{\alpha\in F}|c_\alpha|^2
 \le C_q^2\|\eta\|^2\sum_{\alpha\in F}|c_\alpha|^2 .
\end{equation}
All finite characters outside a used word's sites have image zero; their coefficients may be retained on the right. After combining repeated source group terms, Parseval in the actual cyclic-lamp group basis gives \(\sum|c_\alpha|^2=\|x\|_2^2\). Equation (HS11) is (HS1). Since \(\eta\) was arbitrary with the same bound, taking its unit-sphere supremum proves (HS2).

The scalar Cauchy--Schwarz route has exactly the needed summability: for every \(\zeta\in H\),
\begin{equation}\tag{HS12}
 \sum_\alpha|\langle z_\alpha,\zeta\rangle|^2=\langle S_\eta\zeta,\zeta\rangle\le C_q^2\|\eta\|^2\|\zeta\|^2.
\end{equation}
For finite \(F\) the sum equals \(\langle S_{\eta,F}\zeta,\zeta\rangle\), which is bounded by (HS10). The full equality in (HS12) uses the operator-norm convergence of that positive finite-subset net to \(S_\eta\) proved after (HS9). For any \(c\in\ell^2\) of the source basis, (HS12) implies
\begin{equation}\tag{HS13}
 \sum_\alpha|c_\alpha\langle z_\alpha,\zeta\rangle|
 \le\|c\|_2\left(\sum_\alpha|\langle z_\alpha,\zeta\rangle|^2\right)^{1/2}<\infty.
\end{equation}
This scalar absolute convergence does not assert absolute vector or operator convergence of the ungrouped source-basis series. Finite-tail application of (HS11) gives unconditional vector norm convergence, and the uniform estimate (HS2) gives operator-norm convergence of \(T_0\) applied to \(L^2\) partial sums. Thus its completion is one well-defined bounded linear map, independent of the enumeration or net of finite subsets. No common dense countable subset of \(H\) is required.

A useful equivalent two-stage Cauchy--Schwarz check makes the exponent explicit. For one word and arbitrary unit \(\eta,\zeta\), equations (HS3)--(HS4) give
\[
 \mathbb E|\langle W_g\eta,\zeta\rangle|^2=\langle S_{\eta,g}\zeta,\zeta\rangle\le q^{|g|}.
\]
Probability-space Cauchy--Schwarz, applied to \(a_g\) and this scalar coefficient, therefore gives
\begin{equation}\tag{HS13a}
 \|\mathbb E[a_gW_g]\|_{\mathrm{op}}\le q^{|g|/2}\|a_g\|_2.
\end{equation}
For any \(L^2\) family \((a_g)\), sequence-space Cauchy--Schwarz gives
\begin{equation}\tag{HS13b}
 \begin{aligned}
 \sum_g\|\mathbb E[a_g W_g]\|_{\mathrm{op}}
 &\le\sum_g q^{|g|/2}\|a_g\|_2\\
 &\le \left(\sum_g q^{|g|}\right)^{1/2}\left(\sum_g\|a_g\|_2^2\right)^{1/2}
 = C_q \left(\sum_g\|a_g\|_2^2\right)^{1/2}.
 \end{aligned}
\end{equation}
For general \(L^2\) coefficients \(a_g\), \(\mathbb E[a_g W_g]\) is the finite operator sum
\(\sum_{a\in\Omega_{Z_g}} W_g(a)\int_{\omega|_{Z_g}=a} a_g\,d\mu\).
These scalar integrals exist since \(L^2\) is contained in \(L^1\) on this probability space; (HS13a) follows by \(L^2\) density. Thus the word-grouped coefficient series is absolutely operator-norm convergent. This is not absolute convergence of the ungrouped Fourier rank-one series in (HS8).
It requires \(\sum_g q^{|g|}<\infty\), not \(\sum_g q^{|g|/2}<\infty\); replacing the weighted Cauchy--Schwarz step by the latter requirement would wrongly strengthen the sufficient threshold to \(q<1/9\). Both estimates follow from the same averaged frame identity.

\subsubsection{Identification with the spatial paired norm}

For the standard-form specialization let \(N\) act faithfully normally by \(L\) on \(H\), with conjugation \(J\) and \(r(x)=JL(x^*)J\). Given \(v_j\in U(N)\), set \(U_j=r(v_j^*)=JL(v_j)J\). Define \(V:H\otimes\overline H\to S_2(H)\) by \(V(\eta\otimes\overline\zeta)=\lvert\eta\rangle\langle\zeta\rvert\), and \(t(X)=JX^*J\) on \(S_2(H)\).
Rank-one inner products show that \(V\) and \(t\) are complex-linear unitaries; \(t\) sends \(\lvert\eta\rangle\langle\zeta\rvert\) to \(\lvert J\zeta\rangle\langle J\eta\rvert\) and \(t^2=\mathrm{id}\). Finite-rank density extends these identities even if \(H\) is nonseparable. Directly,
\[
 t(U_j X U_j^*)=L(v_j)t(X)L(v_j)^*,
\]
so \(V^{-1}tP_+t^{-1}V=n^{-1}\sum_j L(v_j)\otimes\overline{L(v_j)}\).
Both spatial representations are faithful, whence
\begin{equation}\tag{HS14}
 q=\|P_+\|_{S_2}=\left\|n^{-1}\sum_jv_j\otimes\overline{v_j}\right\|_{N\otimes_{\min}\overline N}.
\end{equation}
The map \(t\) is not being composed into \(T_0\) as a CP map. This is solely a unitary equivalence of Hilbert-space operators, and uses the conjugate algebra, not an unannounced opposite-algebra identification.

If \(\phi\) is any normal state on \(N\), its natural-cone vector \(\xi_\phi\) has norm \(1\) and can be substituted for \(\eta\) in (HS1). Faithfulness is unnecessary: \(H,L,J\) remain the standard form of all \(N\), independent of \(\phi\). No state-support compression or right multiplication on the potentially defective GNS quotient of \(\phi\) is used. In fact (HS1) holds for all other target vectors at once, so the estimate has no hidden state-dependent constant.
This proves Lemma HS. In the construction of Section~3.1, \(H=H_N\) and \(U_j=r_N(v_j^*)\) give exactly \(W_g=r_N(B_g^*)\), with the same slots. Equation (HS14) identifies its \(q\) with the one in Proposition M. Hence the actual compression satisfies \(\|T_0(x)\|\le C_q\|x\|_{2,D}\), uniformly in every target vector. Its positivity is supplied by Section~3.2, separately from this norm estimate. No matrix-degree-independent \(L^2\) bound is used.

\subsection{The completed reduced kernel and one normal UCP extension}

Let \(p:C_f\to C\) be the regular quotient. For \(x\in\ker p\) choose \(a_m\in A\) converging in full norm to \(x\). Then \(\|a_m\|_2\le\|p(a_m)\|\to0\). Full boundedness of \(\Phi\) and the estimate of Lemma HS give \(\Phi(x)=0\). Thus \(T(p(x))=\Phi(x)\) is one unital contraction \(C\to L_N(N)'\). For positive \(X\in M_k(C)\), lift \(\sqrt X\) to \(Y\in M_k(C_f)\); then \(T^{(k)}(X)=\Phi^{(k)}(Y^*Y)\ge0\). This proves UCP at each \(k\) without assuming a CP section of the quotient.

Here is the complete extension argument. Since the trace on \(D\) is a faithful normal state, \(D\) is contained in \(L^2(D)\); \(A\) is \(L^2\)-dense by the group basis and norm-dense in \(C\). Complete the one \(L^2\)-bounded map to
\[
 \widetilde T:L^2(D)\to B(H_N),\qquad \|\widetilde T\|\le C_q,
\]
and set \(\theta=\widetilde T|_D\). Norm approximation from \(A\) to \(C\) also converges in \(L^2\), so \(\theta\) agrees with \(T\) on \(C\) and preserves the unit. Its values are in the norm-closed algebra \(L_N(N)'\). For \(x\in\operatorname{Ball}(D)\), a Kaplansky net \(c_\alpha\in\operatorname{Ball}(C)\) converges strongly to \(x\) on source \(L^2(D)\), hence \(\|c_\alpha-x\|_2=\|(c_\alpha-x)\widehat1\|\to0\). Therefore \(T(c_\alpha)\) converges to \(\theta(x)\) in operator norm and \(\|\theta(x)\|\le1\).

For fixed \(k\) and \(X\ge0\) in \(M_k(D)\), take positive Kaplansky approximants \(X_\alpha\in M_k(C)\), strongly converging to \(X\). Their entries converge in \(L^2\), and
\[
 \|\theta^{(k)}(X)-T^{(k)}(X_\alpha)\|
 \le C_q\sum_{ij}\|X_{ij}-(X_\alpha)_{ij}\|_2\to0.
\]
Thus \(\theta\) is CP. Once positive entrywise \(L^2\) approximants are given, this last implication does not require an additional uniform norm bound on them. The bounded unit-ball approximation above served to prove contractivity. Matrix \(k\) remains fixed during the finite sum.

If \(0\le x_\alpha\) increases to \(x\) in \(D\), normality of its source trace gives
\[
 \|x-x_\alpha\|_2^2\le\|x\|\tau_D(x-x_\alpha)\to0 .
\]
Consequently \(\theta(x_\alpha)\) increases to \(\theta(x)\), even in operator norm, proving normality. These are nets, with no separability assumed. Uniqueness among normal extensions of \(T\) follows from weak density of \(C\).
The moments (MOM) are unchanged. Normality of the raw \(\pi_f\) on all \(H_{\mathrm{raw}}\), and normality of a later Calkin pullback, have not been inferred.

\subsection{Finite group models and the faithful \texorpdfstring{\(C^*\)}{C*} copy}

First construct the required sequence. For a nontrivial reduced word in \(F(s,t)\), assign distinct vertices \(0,\ldots,m\) to its path, read in action order (rightmost letter first for usual composition of permutations). For a positive letter \(u\) prescribe its permutation to send the current vertex to the next; for a negative letter prescribe the inverse step. For each generator these prescriptions are injective partial bijections: a conflicting domain or range could occur only at an adjacent inverse pair, which reducedness excludes. Complete each partial bijection to a permutation of the finite vertex set. The represented word moves \(0\) to \(m\). Thus every nontrivial word survives a finite quotient. Intersect the finitely many kernels required to separate all words of length \(\le2l\); the resulting quotient \(\pi_l:G\to G_l\) is injective on the ball of radius \(l\).

For a finite-support lamp \(a\) define
\[
 (\pi_{l!}a)(\overline h)=\sum_{\pi_l(h)=\overline h} a(h)\pmod{n^2}.
\]
This respects addition and \(\pi_{l!}(g\cdot a)=\pi_l(g)\cdot(\pi_{l!}a)\), so \((a,g)\mapsto(\pi_{l!}a,\pi_l(g))\) is an exact homomorphism onto \(\Gamma_l=(\mathbb Z/(n^2))^{(G_l)}\rtimes G_l\).
If \(g\ne1\) it survives eventually. If \(g=1\) and \(a\ne0\), injectivity on a ball containing its finite support prevents cancellation of distinct lamp sites, so that element also survives eventually. The regular finite representations \(\rho_l\) of \(\Gamma_l\) have dimension \(d_l=|\Gamma_l|\) and \(\operatorname{tr}_l(\rho_l(\gamma))\) eventually equal to \(\delta_{\gamma,1}\) for each fixed \(\gamma\). They therefore give the required trace-convergent sequence. Relations, inverses and repeated words are preserved exactly at every \(l\); the target cocycle need not factor through these finite quotients.

Use any models as in (M). \(I\) is a norm-closed two-sided \(*\)-ideal in \(P\), by the tracial inequalities \(\|ab\|_2,\|ba\|_2\le\|a\|\|b\|_2\) and uniform operator bounds. The trace \(\tau_P(b)=\lim_\omega\operatorname{tr}_l(b_l)\) is a tracial state of \(P\). Cauchy--Schwarz gives \(\tau_P(i)=0\) for \(i\in I\), so it descends to a trace \(\tau_U\) on \(U=P/I\). It is faithful: \(\tau_U(x^*x)=0\) says exactly that any representative of \(x\) belongs to \(I\).
Let \(\beta:C_f\to P\) be \(\beta(x)=(\rho_l(x))_l\), and \(q_\tau:P\to U\) the quotient. Explicitly \(S_u=\beta(w_u)\), \(P_{u,j}=\beta(p_{1,u,j})\), and \(b_{u,j}=S_uP_{u,j}\). Since the \(p_{1,u,j}\) are exact finite Fourier polynomials forming orthogonal partitions, these are the exact model unitaries and projections used in Proposition M.
The trace limits give \(\tau_Uq_\tau\beta=\tau_Dp\) on \(C_f\) by norm density.
The regular trace on \(C\) is faithful: a positive element with trace \(0\) has square root killing \(\delta_1\), hence every basis vector by the commuting right regular action. Apply the trace identity to \(x^*x\) to obtain
\[
 \ker(q_\tau\beta)=\ker p .
\]
There is therefore a faithful unital \(*\)-homomorphism
\begin{equation}\tag{COPY}
 \sigma_C:C\to U,\quad \sigma_Cp=q_\tau\beta,\quad
 \sigma_C(d_{u,j})=q_\tau(b_{u,j}).
\end{equation}
This establishes the \(C^*\) norm, not just an \(L^2\) identity. It does not infer operator norm convergence of finite blocks from trace convergence.

\subsection{A trace-GNS return and a map on the whole Calkin algebra}

In the trace GNS spaces define the densely specified map \(W:L^2(D)\to L^2(U,\tau_U)\) on the dense \(C\) by
\[
 W(\widehat c)=\widehat{\sigma_C(c)}.
\]
Trace preservation makes \(W\) an isometry on the completion, with norm \(1\) on all \(L^2(D)\). Right multiplication \(R_U(a)\) on \(L^2(U,\tau_U)\) is bounded by \(\|a\|\) and has adjoint \(R_U(a^*)\); the tracial property gives these assertions even though \(U\) is only a \(C^*\) algebra. For \(c\in C\),
\[
 WR_D(c)=R_U(\sigma_C(c))W,\qquad WR_D(c^*)=R_U(\sigma_C(c^*))W .
\]
Thus \(W\) intertwines both right multiplication and its adjoint. For \(x\in U\), the compression \(W^*L_U(x)W\) commutes with \(R_D(C)\), hence \(R_D(D)\), since \(L_U(x)\) commutes with right multiplication and \(\operatorname{range}(W)\) reduces \(R_U(\sigma_C(C))\). The standard commutant theorem gives a unique element \(F_U(x)\in D\) with
\[
 L_D(F_U(x))=W^*L_U(x)W .
\]
Compression at every matrix level proves UCP. The same dense-vector calculation gives \(L_U(\sigma_C(c))W=WL_D(c)\), so \(F_U(\sigma_C(c))=c\) for \(c\in C\).
No assertion that \(U\) is a von Neumann algebra, or that a normal embedding \(D\to U\) has been constructed, is used.

Let \(q_c:B(H_c)\to Q_c\) be the Calkin quotient and let \(\Delta:B(H_c)\to P\) be diagonal block compression. It is UCP and sends compacts into the norm-null block ideal \(c_0\subset I\): a rank-one block norm is \(\|\eta_l\|\|\zeta_l\|\to0\), and finite-rank norm approximation finishes. Consequently \(F_Uq_\tau\Delta\) kills \(K(H_c)\) and defines a UCP map on all \(Q_c\), by positive matrix quotient lifts:
\begin{equation}\tag{ALL-Q}
 \begin{aligned}
 F_Q(q_c(X))&=F_U(q_\tau(\Delta(X))),\\
 F_Q(j(b))&=F_U(q_\tau(b))\qquad\text{for every }b\in P.
 \end{aligned}
\end{equation}
Put \(T_v=\theta F_Q\). (COPY), (MOM) and (ALL-Q) give simultaneously every specified moment and \(T_v(j(I))=0\), proving (M). The same choices of models, \(\omega,W,F_U,F_Q\) and \(\theta\) serve all moments and every later cut.

Only range containment is needed for \(\theta\) and \(T_v\). Source surjectivity of a return, when separately proved, would not imply target surjectivity. Nor is a prescribed reverse assignment \(d_{u,j}\mapsto j(b_{u,j})\) assumed positive. The faithful copy is in \(P/I\); no UCP lift of it to \(P\) or \(Q_c\) is an input. Normality of \(\theta\) does not make the pullback of \(T_v\) to \(B(H_c)\) normal. These distinctions prevent strengthening (M) beyond its proof.

\section{Transfer to the original \texorpdfstring{\(C^*\)}{C*}-algebra}

\paragraph{Proposition D (proposed detector). Assumptions and conclusion.}
Let \(B\) be a nonzero exact \(C^*\)-algebra, possibly nonunital and nonseparable,
\(V=B^{**}\), and \(\kappa:B\to V\) the canonical inclusion. Assume G has supplied
a nonzero \(e\in\operatorname{Proj}(V)\), \(N=eVe\) in normal standard form
\(L_N,r_N\) on \(H_N\), relative unitaries \(v_1,\ldots,v_n\), and a normal state
\(\phi\) with unit cone vector \(\xi\) such that
\[
 r=\frac1n\sum_j \operatorname{Re}\langle L_N(v_j)r_N(v_j^*)\xi,\xi\rangle>\frac{63}{64}.
\]
The small-\(q\) bound used to obtain M is part of that supplied input; D
will not assert any small-\(q\) bound for its replacement coefficients.

Assume the following fixed M data, including its faithful \(C^*\) source model.
In this section \(T\) denotes the same \(T_v\) and \(\sigma\) denotes \(\sigma_C\) of M.
\(G=F(s,t)\), \(\Gamma=(\mathbb Z/(n^2))^{(G)}\rtimes G\), \(C=C_r^*(\Gamma)\).
The original cyclic lamp at each site has its uniform \(n\)-by-\(n\) spectral
table; \(p_{h,u,j}\), \(u=s,t\), are row/column projections and
\(\sigma_g(p_{h,u,j})=p_{gh,u,j}\). Write \(d_{u,j}=w_u p_{1,u,j}\) in \(C\).
Take \(P=\prod_{\ell\geq1} M_{d_\ell}\), with positive finite integers \(d_\ell\) and a
free ultrafilter \(\omega\), and \(I=\{b\in P:\lim_\omega\operatorname{tr}_\ell(b_\ell^*b_\ell)=0\}\).
Write \(U=P/I\), quotient \(q_\tau:P\to U\), \(H_c=\bigoplus_\ell\mathbb C^{d_\ell}\),
\(Q_c=\mathcal B(H_c)/\mathcal K(H_c)\), and the block quotient \(*\)-map \(j:P\to Q_c\).
In \(P\) let \(S_u\) be unitaries and let \(P_{u,1},\ldots,P_{u,n}\) be orthogonal
projections summing to \(1\), with \(b_{u,j}=S_u P_{u,j}\). Assume
\(\sigma:C\to U\) is a faithful \(*\)-embedding with \(\sigma(d_{u,j})=q_\tau(b_{u,j})\),
and \(T:Q_c\to L_N(N)'\) is one UCP map satisfying
\begin{equation}
 T(j(b_{u,j}))=\frac1n r_N(v_j^*),\qquad T(j(i))=0\quad(i\in I).
 \tag{M-IN}
\end{equation}
All these maps and the free ultrafilter are fixed before any coefficient
in \(B\) or ideal cut is chosen. \(j(I)\) need not vanish in \(Q_c\). No global
Stinespring representation is assumed to kill \(j(I)\). Normality of the
lift \(T\circ(\mathcal B(H_c)\to Q_c)\) and a normal von Neumann embedding of
the source are not assumed.

The target is one state \(\Psi\) on \(E=B\otimes_{\max} Q_c\) and a selfadjoint
finite tensor \(u\in B\odot Q_c\) with
\[
 \|u\|_{\max}\leq1,\qquad \|u\|_{\min}<7/8,\qquad \Psi(u)>251/256.
\]
It follows that \(k=(u-7/8)_+\) lies in \(E\), has zero canonical min image and
\(\Psi(k)>27/256\). All left coefficients remain in \(B\), even when it lacks a unit.

\subsection{One fixed state and its normal left variable}

Set \(\alpha:V\to\mathcal B(H_N)\), \(\alpha(X)=L_N(eXe)\). This is normal UCP, with
\(\alpha(1_V)=1_{H_N}\), and its range commutes with \(T(Q_c)\). It is generally
not a \(*\)-representation. We construct the needed commuting dilation
while retaining the compression as a UCP map.
On \(V\odot Q_c\odot H_N\) use the sesquilinear form, linear in its first variable,
\[
 \langle X\otimes z\otimes\eta,Y\otimes w\otimes\zeta\rangle
 =\langle\alpha(Y^*X)T(w^*z)\eta,\zeta\rangle.
\]
For finite families \(A=[\alpha(X_i^*X_j)]\) and \(D=[T(z_i^*z_j)]\) are positive
matrices and each entry of \(A\) commutes with each entry of \(D\). Choose
\(S=A^{1/2}\) in the matrix \(C^*\)-algebra generated by \(\alpha(V)\). Functional
calculus stays in that matrix algebra, so every entry of \(S\) commutes
with every entry of \(D\). This choice is essential; an arbitrary \(S\) with
\(S^*S=A\) need not have that property. With \(D_\ell=\operatorname{diag}_i(S_{\ell i})\),
\[
 [A_{ij}D_{ij}]=\sum_\ell D_\ell^* D D_\ell\geq0.
\]
This proves positivity, retaining all cross terms; it does not assume
entries within \(A\) commute among themselves. Applying the same argument
to \([\alpha(X_i^*(\|X\|^2-X^*X)X_j)]\) proves that left multiplication by \(X\)
on the first factor is bounded by \(\|X\|\). The corresponding positive
matrix for \(T\) bounds multiplication on the second factor by \(\|z\|\).
Quotienting the null space and completing gives commuting unital
\(*\)-representations \(\rho:V\to\mathcal B(K)\), \(\nu:Q_c\to\mathcal B(K)\), and the isometry
\(J\eta=[1\otimes1\otimes\eta]\), with
\begin{equation}
 J^*\rho(X)\nu(z)J=\alpha(X)T(z). \tag{COMP}
\end{equation}
Thus \(J\xi\) defines a state \(\Omega_V\) on \(V\otimes_{\max} Q_c\). No assertion about
minimal continuity of this product representation is used.

The left representation is normal: on finite sums of the dense vectors
\(\rho(X_i)\nu(z_i)J\eta_i\), all \(\rho(A)\)-coefficients are finite sums of
\[
 \langle\alpha(X_j^* A X_i)T(z_j^* z_i)\eta_i,\eta_j\rangle,
\]
which are normal in \(A\).
Approximation of arbitrary vectors in norm, using \(\|\rho(A)\|\leq\|A\|\),
gives normal coefficients because \(V_*\) is norm closed in \(V^*\). This uses
normal \(\alpha\) and fixed \(T\)-values, never normality in the \(z\) variable.

Pull \(\Omega_V\) back through \(\kappa\otimes_{\max}\mathrm{id}\) to original \(B\), obtaining
\begin{equation}
 \Psi(b\otimes z)=\langle L_N(e\kappa(b)e)T(z)\xi,\xi\rangle. \tag{STATE}
\end{equation}
The pullback map need not be injective. The restriction is a state,
not merely a substate: for a positive contractive approximate identity
\(b_a\) in \(B\), normality gives \(\Psi(b_a\otimes1)\to\phi(e)=1\). Hence its norm is \(1\).
There is one fixed vector state and commuting dilation throughout.
Only the left \(V\)-variable is called normal; no predual/normality is
being assigned to the \(C^*\)-algebra \(Q_c\) or to the whole maximal \(C^*\) product.

\subsection{Finite normal tests and original-\texorpdfstring{\(B\)}{B} coefficients}

For \(k=1,\ldots,n\) define the normal slice in the same state
\[
 f_k(X)=\frac1n\langle L_N(eXe)r_N(v_k^*)\xi,\xi\rangle,\qquad \|f_k\|\leq1/n.
\]
Equivalently it is the \(\Omega_V\)-slice at \(j(b_{u,k})\) for either \(u\). Its
normality follows from \(L_N\) and compression, since \(r_N(v_k^*)\xi\) is a
fixed vector. In canonical \(V=B^{**}\), these normal functionals belong to
the specified predual \(B^*\), so Goldstine applies to them.

After fixing \(\Psi,T\), all \(v_j\) and all these tests, Goldstine says that
\(\kappa(\operatorname{Ball}(B))\) is \(\sigma(B^{**},B^*)\)-dense in \(\operatorname{Ball}(B^{**})\). Each \(v_j\) has norm \(1\)
in \(V\) even though its unit is \(e\). Therefore choose \(a_j\in\operatorname{Ball}(B)\) with
\begin{equation}
 |f_k(\kappa(a_j)-v_j)|<\gamma=\frac1{256n}\quad\text{for every \(k\) and \(j\)}. \tag{GOLD}
\end{equation}
Each \(j\) is an elementwise choice against one finite common list. For
example, the inequalities define a weak*-open neighborhood of \(v_j\);
Goldstine puts a point of \(\kappa(\operatorname{Ball}(B))\) in it. There are only finitely
many \(j\), so all choices can be made before forming the tensor below.
No sequence or metrizability is assumed. The chosen elements belong
to \(B\) itself, not only its unitization. They need not be unitaries,
supported on \(e\), multiplicative images of \(v_j\), or a completely contractive
linear interpolation on their span. No products of the \(a_j\) are being
approximated, and no matrix-level or UCP-map approximation is claimed.
Kaplansky density and local reflexivity are not used in this selection.

Using the same coefficients set
\begin{align*}
 y&=\tfrac14\sum_{u=s,t;\,j}(a_j\otimes d_{u,j}+a_j^*\otimes d_{u,j}^*),\\
 t_0&=\tfrac14\sum_{u,j}(a_j\otimes b_{u,j}+a_j^*\otimes b_{u,j}^*),\\
 u_0&=(\mathrm{id}_B\otimes j)(t_0),
\end{align*}
in \(B\odot C\), \(B\odot P\) and \(B\odot Q_c\), respectively.
The expression with \(v_j\) instead of \(\kappa(a_j)\) in \(\Omega_V\) has value \(r\).
For each of the \(2n\) nonadjoint terms, (M-IN) and (GOLD) bound its state-value
error by \(\gamma\). Adjoint paired terms have conjugate values because
\(\Omega_V\) and \(\Psi\) are positive states. Thus no separate assertion that
the \(f_k\)-values of \(a_j^*\) approximate \(v_j^*\) is required. We get
\begin{equation}
 |\Psi(u_0)-r|<\tfrac14(2)(2n)\gamma=1/256,\qquad \Psi(u_0)>251/256. \tag{LOWER}
\end{equation}
The approximation has preserved specified values of one fixed state,
not the entire original joint law or a small paired tensor norm for \(a_j\).

\subsection{An independent spatial upper bound}

Take a faithful nondegenerate representation \(\pi_B\) of original \(B\) and
the faithful regular crossed-product representation of \(C\) on the fibers
\(\ell^2(G,L^2(\Omega)\otimes H_B)\). On its \(h\) fiber define
\[
 C_u(h)=\sum_j L(\sigma_{h^{-1}}(p_{1,u,j}))\otimes\pi_B(a_j).
\]
The lamp projections are orthogonal and commute with the coefficient
operators, so every \(C_u(h)\) is a contraction. With \(\lambda_u\) the left
shift on \(G\), \(Z_u=\lambda_u C_u\) represents \(\sum_j a_j\otimes d_{u,j}\).
At a fiber \(h\), \(Z_u\xi\) has norm at most \(\|\xi(u^{-1}h)\|\), and \(Z_u^*\xi\) has
norm at most \(\|\xi(uh)\|\). Hence the norm of \(\sum_u(Z_u+Z_u^*)\) is bounded
by the scalar adjacency norm of the four-regular tree.
For completeness, its weight \(w(h)=3^{-|h|/2}\) gives row weighted-sum
ratio \(2\sqrt3\) off the root and \(4/\sqrt3\) at the root. Weighted
Cauchy--Schwarz gives
\[
 |Af(h)|^2\leq2\sqrt3\sum_{k\sim h}\frac{w(h)}{w(k)}|f(k)|^2;
\]
summing \(h\) and applying the same ratio bound gives
\(\|A\|\leq2\sqrt3\). The positive weight need not itself lie in \(\ell^2\).
Thus in these faithful spatial representations
\begin{equation}
 \|y\|_{B\otimes_{\min} C}\leq\sqrt3/2. \tag{SOURCE-MIN}
\end{equation}
This holds for every tuple \(a_j\in\operatorname{Ball}(B)\), irrespective of its quality
of approximation to \(v_j\). It is not weak*-continuity of a tensor norm.

Faithfulness of \(\sigma\) and injectivity of the min tensor norm give
\begin{equation}
 s=\|(\mathrm{id}_B\otimes q_\tau)(t_0)\|_{B\otimes_{\min} U}
  =\|(\mathrm{id}_B\otimes \sigma)(y)\|=\|y\|_{\min}\leq\sqrt3/2. \tag{QUOT}
\end{equation}
No small min norm for uncut \(u_0\) in the Calkin algebra is asserted.
In contrast, any commuting representations \(\pi_B,\nu_P\) of \(B\) and \(P\) give
\[
 \Big\|\sum_j\pi_B(a_j)\nu_P(P_{u,j})\Big\|\leq1
\]
by orthogonal blocks. Multiplying
by \(\nu_P(S_u)\) preserves the bound, so averaging two contractions and
their adjoints gives
\begin{equation}
 \|t_0\|_{B\otimes_{\max}P}\leq1,\qquad \|u_0\|_{\max}\leq1. \tag{MAX}
\end{equation}
This uses no injectivity of an arbitrary maximal tensor inclusion.
\subsection{A finite ideal cut preserving that same state}

Set \(\eta=1/256\). This is the only use of exactness: for the norm-closed
two-sided ideal \(I\) of \(P\),
\[
 \ker(B\otimes_{\min} P\to B\otimes_{\min} U)=\overline{B\odot I}.
\]
The quotient norm and density therefore supply a finite
\(K_0=\sum_i c_i\otimes z_i\in B\odot I\) with \(\|t_0-K_0\|_{\min}<s+\eta/2\).
To see the finite choice, first approximate the quotient norm by an
element of the closed kernel and then approximate it by \(B\odot I\).
This is not an appeal to local reflexivity or to a completely bounded
section of \(P\to U\). In particular every \(c_i\) is in original \(B\).

Let \(A_0=\sum_i\|c_i\|\). If \(A_0=0\), \(K_0=0\) and take \(h=0,c=1\). Otherwise put
\begin{align*}
 Z&=\sum_i(z_i^*z_i+z_i z_i^*)\in I_+,\qquad \delta=(\eta/A_0)^2,\\
 h&=Z(Z+\delta)^{-1}\in I,\qquad c=1-h\in P.
\end{align*}
The inverse uses the unit of \(P\), not a unit adjoined to \(B\). Here \(0\leq c\leq1\),
\(q_\tau(c)=1\) and \(\|c Z c\|\leq\delta/4\) by scalar functional calculus. Since
\(z_i^*z_i\) and \(z_i z_i^*\) are bounded by \(Z\), this yields
\[
 \|z_i c\|,\|c z_i\|\leq\sqrt\delta/2,\qquad
 \|c K_0 c\|_{\min}\leq A_0\sqrt\delta/2=\eta/2.
\]
The notation \(cK_0c\) means multiplication only in the second factor.
It is bounded also for nonunital \(B\), as can be checked in every spatial
or commuting representation; it does not move the first coefficients
to a unitization. Define the finite selfadjoint tensors
\[
 t=(1\otimes c)t_0(1\otimes c)\in B\odot P,\qquad
 u=(\mathrm{id}_B\otimes j)(t)\in B\odot Q_c.
\]
Then \(\|u\|_{\max}\leq1\) by (MAX) and contractive multiplication, and
\begin{equation}
 \|u\|_{\min}\leq\|t\|_{\min}<s+\eta\leq\sqrt3/2+1/256<7/8. \tag{CUT}
\end{equation}
The last strict margin follows from \((223/256)^2-3/4=577/65536>0\).

For every \(b\in P\), \(cbc-b=-hb-bh+hbh\in I\). The fixed (M-IN) identity
therefore gives \(T(j(cbc))=T(j(b))\) before or after the finite cut was
chosen. Applying (STATE) term by term gives
\begin{equation}
 \Psi(u)=\Psi(u_0)>251/256 \quad\text{exactly}. \tag{SAME}
\end{equation}
No new state extension is selected. No multiplicativity of \(T\) is used;
no convergence argument is made in \(Q_c\) or through its lift. The
ideal elements involved lie in \(P\), and only the values on \(j(P)\) enter.
There is no assertion that a whole-\(Q_c\) representation annihilates \(j(I)\).

\subsection{Functional calculus and nonunital \texorpdfstring{\(B\)}{B}}

In \(E=B\otimes_{\max} Q_c\) define \(k=(u-7/8)_+\) by continuous functional
calculus. If \(E\) is nonunital, do the calculation in its unitization;
\(f(t)=(t-7/8)_+\) has \(f(0)=0\), so \(k\) belongs to \(E\) itself. For the canonical
\(*\)-homomorphism \(\chi:E\to B\otimes_{\min} Q_c\), functional calculus and (CUT)
give \(\chi(k)=0\) exactly. The scalar inequality \((t-7/8)_+\geq t-7/8\) gives
\(k\geq u-(7/8)1\) in the unitization of \(E\). Extend the state unitally to obtain
\[
 \Psi(k)\geq\Psi(u)-7/8>27/256,\qquad 0\leq k,\quad \|k\|_{\max}\leq1/8.
\]
No functional calculus through a UCP map is asserted. Only \(u\), not \(k\),
is necessarily an algebraic finite tensor. \(H_c\) is a countably infinite
sum of nonzero finite blocks, so one Hilbert-space unitary identifies
\(Q_c\) with the fixed separable Calkin algebra, leaving original \(B\) intact.

\subsection{Scope and order of the approximations}

Goldstine selects finitely many elements \(a_j\) of \(\operatorname{Ball}(B)\), testing only
normal functionals in the canonical \(V\) variable. It approximates neither
matrices nor a UCP map nor their multiplication. (SOURCE-MIN) supplies
an independent universal contraction bound. Exactness subsequently
selects a finite element of \(B\odot I\) for one quotient-norm approximation.
The fixed state is preserved by an exact all-\(I\) identity, not by a
weak* limit of its right-hand variable. No faithful normal state,
support cut, local-reflexivity hypothesis, maximal-inclusion isometry,
or equality of the min and state-evaluating representations is added.
D is a conditional transfer result under exactly the displayed G/M
package and exactness of \(B\). Propositions G and M supply that package in
Sections 2--3, and Section 1 derives exactness from the original hypothesis.
This proves Proposition D.

\section{Assembly of the proposed theorem}

Assume the canonical \(\chi_B\) is injective for a nonzero \(B\). Section 1 gives
exactness of this original \(B\). If \(B\) were nonnuclear, input (B) would make
its canonical \(V=B^{**}\) noninjective. Proposition G, with \(q_0=1/4\) and
\(\epsilon=1/32\), supplies \(e\) in \(V\), \(N=eVe\), one indexed unitary tuple \(v\) and
one normal state \(\phi\) with cone vector \(\xi\) and \(q\leq1/4,r>63/64\). The
type III amplification is absorbed normally and onto before \(e\) and \(\phi\)
are chosen; no left matrix algebra remains.

Use this same \(N\) and tuple in Proposition M. Fix one trace-convergent
finite source sequence and one free ultrafilter after \(n\) is fixed.
Lemma HS has \(C_q^2\leq5\). Section 3.4 therefore constructs one normal UCP
\(\theta\) on the finite source \(D\); no target trace or faithful \(\phi\) is needed.
Sections 3.5--3.6 construct exact partitions, faithful \(\sigma_C\) and one
whole-\(Q_c\) moment map \(T_v\). All specified moments and all-\(I\) identities
hold simultaneously, before any original-\(B\) coefficient or ideal cut.

Proposition D now fixes \(\alpha\) and \(\Psi\) using that same \(e,\phi,\xi,T_v\).
Only then are \(a_j\) chosen in original \(\operatorname{Ball}(B)\). Their spatial upper bound
uses the regular source representation and holds for every contractive
tuple, whereas their lower bound is measured in the fixed state \(\Psi\).
The faithful copy identifies the quotient norm of that same tensor \(t_0\).
Exactness of \(B\) selects one finite ideal correction \(K_0\) and \(c=1-h\). The
identity \(cbc-b\in I\), valid for every \(b\in P\), preserves \(\Psi\) exactly.
Functional calculus in \(B\otimes_{\max} Q_c\) gives \(k\geq0\) with
\[
 \chi(k)=0,\qquad \Psi(k)>27/256,\qquad \|k\|_{\max}\leq1/8.
\]

The block Hilbert space \(H_c\) is a countably infinite sum of nonzero
finite-dimensional spaces. Choose a unitary \(H_c\to H_*\). Conjugation
identifies \(Q_c\) with \(Q_*\) and intertwines both tensor completions and
their canonical map. It sends \(k\) to a nonzero positive element of
\(\ker\chi_B\), a contradiction. Thus \(B\) is nuclear. Together with the
immediate nuclear direction this completes the proposed proof of (T).

\paragraph{Dependency table.} The following are proved steps, not additional axioms.

\begin{center}
\small
\begin{tabular}{@{}p{0.34\linewidth}p{0.22\linewidth}p{0.35\linewidth}@{}}
\hline
Step & Proof location & Standard inputs\\ \hline
Hypothesis \(\Rightarrow\) exact \(B\) & Section 1 & (E), maximal ideal row\\
Literal \(\sigma\)-finite corner & Section 2.1 & Injectivity/compactness\\
Finite retract in type III & Sections 2.2--2.3 & (C), proper infiniteness\\
Small \(q\) and one normal \(\phi\) & Sections 2.4--2.5 & (H), (S), separation\\
Full compression/positivity & Sections 3.1--3.2 & Full group \(C^*\) completion\\
Uniform operator \(L^2\) bound & Section 3.3 & Spatial min, HS theory\\
One normal source UCP & Section 3.4 & Kaplansky, finite trace\\
Finite models/faithful copy & Section 3.5 & \(C^*\) GNS and quotient\\
All-Calkin moment map & Section 3.6 & Source tracial commutant\\
One common tensor state & Section 4.1 & Commuting representations\\
Coefficients in original \(B\) & Section 4.2 & Goldstine in \(B^{**}\)\\
Same-tensor min/max bounds & Section 4.3 & Spatial/commuting norms\\
Exact cut and fixed state & Section 4.4 & Exactness of original \(B\)\\
Zero min image, nonzero \(k\) & Section 4.5 & \(C^*\) functional calculus\\ \hline
\end{tabular}
\end{center}

G uses no M or D. M requires only the displayed tuple and \(q\), not the
target state or \(B\). D uses the already fixed G/M outputs. No hypothesis
is asserted for \(N\), a unitization of \(B\) or an amplification of \(B\). Normality
is used on the source \(D\) for \(\theta\) and on \(V\) for \(\alpha\) and its finite
slices, not on the whole Calkin map. Only block elements \(j(P)\) enter the
cut. The nonseparable standard-form \(H_N\) is never identified with \(H_*\).

\appendix
\section{The scope of block-ideal annihilation}

For this infinite source, trace convergence forces \(d_\ell\to\infty\):
for any \(m\) distinct group elements, their model matrices have normalized
Hilbert--Schmidt Gram tending to \(I_m\), hence rank \(m\leq d_\ell^2\) eventually.
Adding one trivial representation to each model therefore preserves the
trace limit, since its weight is \(1/(d_\ell+1)\). The same trace-kernel proof
then supplies a faithful \(\sigma_C\) for the augmented models. The added
block projection \(p\) belongs to the new trace ideal \(I\) but is not compact:
it fixes infinitely many orthonormal added vectors and has essential
norm \(1\). Applying Section 3.6 for these very models supplies the same simultaneous
moment/all-\(I\) statement used by Section 4.4. No extra dimension assumption is needed.

In any Stinespring representation \(T_v(a)=W^*\nu(a)W\), the guaranteed annihilation subspace is
\[
 K_{\mathrm{block}}=\overline{\operatorname{span}\nu(j(P))W H_N}.
\]
For \(i\in I\) and \(b\in P\), \(ib\in I\) and \(T_v(j((ib)^*(ib)))=0\) give
\(\nu(j(i))\nu(j(b))W=0\). With a commuting left action \(\rho\) this also holds
on \(\overline{\operatorname{span}\rho(V)\nu(j(P))W H_N}\). This is not the whole \(Q_c\)-generated
space. Indeed choose an isometry \(S\) of \(H_c\) onto the infinite added
subspace, \(z=q_c(S)\), \(\bar p=j(p)\). Then
\[
 z^*z=1,\qquad zz^*=\bar p,\qquad z^*\bar p z=1,\qquad
 T_v(\bar p)=0,\quad\text{but}\quad T_v(z^*\bar p z)=1.
\]
Thus full-dilation annihilation is false and unused. Section 4.4 needs
only \(cbc-b\in I\) for \(b\in P\). It uses its actual exactness-selected cut;
replacing that cut by the diagnostic \(p\) alone would not establish
\(\|u\|_{\min}<7/8\).
This control makes the logical boundary explicit, without assuming
multiplicativity, normality or surjectivity of the whole-\(Q\) return.

\begin{thebibliography}{9}
\bibitem{kirchberg}
E. Kirchberg, \emph{The Fubini theorem for exact \(C^*\)-algebras},
Journal of Operator Theory \textbf{10} (1983), 3--8.
Theorem 1.1(ii), (iii), p.~4; countable Hilbert-space convention, p.~3.
\url{https://jot.theta.ro/jot/archive/1983-010-001/1983-010-001-001.pdf}
\bibitem{choieffros}
M. D. Choi and E. G. Effros, \emph{Nuclear \(C^*\)-algebras and injectivity:
the general case}, Indiana University Mathematics Journal \textbf{26} (1977),
443--446. \href{https://doi.org/10.1512/iumj.1977.26.26034}{doi:10.1512/iumj.1977.26.26034}.
Adopted standard theorem.
\bibitem{haagerupcb}
U. Haagerup, \emph{Injectivity and decomposition of completely bounded maps},
Lecture Notes in Mathematics \textbf{1132} (1985), 170--222.
Lemma 2.2, p.~186; conjugate-algebra convention, pp.~185--186.
\bibitem{haagerupstandard}
U. Haagerup, \emph{The standard form of von Neumann algebras},
Mathematica Scandinavica \textbf{37} (1975), 271--283.
Theorem 1.6, p.~275, and Lemma 2.10, pp.~278--279.
\href{https://doi.org/10.7146/math.scand.a-11606}{doi:10.7146/math.scand.a-11606}.
\bibitem{takesaki}
M. Takesaki, \emph{Duality for crossed products and the structure of von Neumann algebras of type III},
Acta Mathematica \textbf{131} (1973), 249--310.
The double-duality theorem is adopted with the scope stated in (C).
\bibitem{hjx}
U. Haagerup, M. Junge and Q. Xu, \emph{A reduction method for noncommutative
\(L^p\)-spaces and applications}, \href{https://arxiv.org/abs/0806.3635v2}{arXiv:0806.3635v2} (2009).
Section 1.1, p.~3, and Section 1.2, p.~5, immediately before (1.10):
the modular crossed product has a normal semifinite faithful trace.
\end{thebibliography}

\end{document}
